\subsection{Requisitos operacionales y experiencia de usuario}
\label{sec:operational-results}

\begin{table}[H]
\centering
\caption{Comparación de requisitos operacionales.}
\label{tab:operational-comparison}
\small
\begin{tabularx}{\textwidth}{lCCCCCY}
\toprule
\textbf{Protocolo} &
\textbf{Registro} &
\textbf{Estado} &
\textbf{Servidor online} &
\textbf{Secreto cliente} &
\textbf{Verif. offline} &
\textbf{Observaciones} \\
\midrule

ECDSA      & \todoresult{} & \todoresult{} & \todoresult{} & \todoresult{} & \todoresult{} & \todoresult{} \\
LWE        & \todoresult{} & \todoresult{} & \todoresult{} & \todoresult{} & \todoresult{} & \todoresult{} \\
Binary-LWE & \todoresult{} & \todoresult{} & \todoresult{} & \todoresult{} & \todoresult{} & \todoresult{} \\
Ring-LWE   & \todoresult{} & \todoresult{} & \todoresult{} & \todoresult{} & \todoresult{} & \todoresult{} \\
LWR        & \todoresult{} & \todoresult{} & \todoresult{} & \todoresult{} & \todoresult{} & \todoresult{} \\
LWE propio & \todoresult{} & \todoresult{} & \todoresult{} & \todoresult{} & \todoresult{} & \todoresult{} \\

\bottomrule
\end{tabularx}
\end{table}

Desde la perspectiva del usuario final, \todoresult{describir cuáles diferencias son
visibles para el usuario y cuáles quedan ocultas por la aplicación}.


% ============================================================
\subsection{Dependencias de confianza y arquitectura}
\label{sec:trust-model}

En lugar de tratar la centralización como una propiedad binaria, se analizan
explícitamente las dependencias de confianza de cada construcción.

\begin{table}[H]
\centering
\caption{Dependencias de confianza e infraestructura.}
\label{tab:trust-model}
\small
\begin{tabularx}{\textwidth}{lCCCCYY}
\toprule
\textbf{Protocolo} &
\textbf{Autoridad} &
\textbf{PKI} &
\textbf{Estado servidor} &
\textbf{Secreto servidor} &
\textbf{Compromiso servidor} &
\textbf{Observaciones} \\
\midrule

ECDSA      & \todoresult{} & \todoresult{} & \todoresult{} & \todoresult{} & \todoresult{} & \todoresult{} \\
LWE        & \todoresult{} & \todoresult{} & \todoresult{} & \todoresult{} & \todoresult{} & \todoresult{} \\
Binary-LWE & \todoresult{} & \todoresult{} & \todoresult{} & \todoresult{} & \todoresult{} & \todoresult{} \\
Ring-LWE   & \todoresult{} & \todoresult{} & \todoresult{} & \todoresult{} & \todoresult{} & \todoresult{} \\
LWR        & \todoresult{} & \todoresult{} & \todoresult{} & \todoresult{} & \todoresult{} & \todoresult{} \\
LWE propio & \todoresult{} & \todoresult{} & \todoresult{} & \todoresult{} & \todoresult{} & \todoresult{} \\

\bottomrule
\end{tabularx}
\end{table}

% ============================================================
\subsection{Comparación global}
\label{sec:global-comparison}

La Tabla~\ref{tab:global-summary} resume las principales métricas cuantitativas.
Esta tabla no pretende producir un ranking absoluto entre protocolos, sino permitir
identificar visualmente los compromisos existentes entre las diferentes alternativas.

\begin{table}[H]
\centering
\caption{Resumen cuantitativo de los resultados principales.}
\label{tab:global-summary}
\resizebox{\textwidth}{!}{
\begin{tabular}{lrrrrrr}
\toprule

\textbf{Protocolo} &
\textbf{Auth. (ms)} &
\textbf{RAM pico} &
\textbf{Comunicación} &
\textbf{Clave pública} &
\textbf{Clave privada} &
\textbf{Seguridad estimada} \\

&
&
\textbf{KB} &
\textbf{bytes} &
\textbf{bytes} &
\textbf{bytes} &
\textbf{bits} \\

\midrule

ECDSA
& \todoresult{}
& \todoresult{}
& \todoresult{}
& \todoresult{}
& \todoresult{}
& \todoresult{} \\

LWE
& \todoresult{}
& \todoresult{}
& \todoresult{}
& \todoresult{}
& \todoresult{}
& \todoresult{} \\

Binary-LWE
& \todoresult{}
& \todoresult{}
& \todoresult{}
& \todoresult{}
& \todoresult{}
& \todoresult{} \\

Ring-LWE
& \todoresult{}
& \todoresult{}
& \todoresult{}
& \todoresult{}
& \todoresult{}
& \todoresult{} \\

LWR
& \todoresult{}
& \todoresult{}
& \todoresult{}
& \todoresult{}
& \todoresult{}
& \todoresult{} \\

LWE propio
& \todoresult{}
& \todoresult{}
& \todoresult{}
& \todoresult{}
& \todoresult{}
& \todoresult{} \\

\bottomrule
\end{tabular}
}
\end{table}


% ============================================================
\subsection{Análisis multidimensional}
\label{sec:multidimensional-analysis}

Una única métrica no permite determinar la conveniencia de una construcción
criptográfica. Un protocolo puede reducir el tiempo de ejecución a costa de incrementar
el tamaño de las claves, mientras que otro puede reducir la comunicación aumentando
la complejidad computacional.

Por esta razón, los resultados se analizan como un conjunto de compromisos entre:

\begin{equation}
    (\text{seguridad},
     \text{tiempo},
     \text{memoria},
     \text{comunicación},
     \text{complejidad}).
\end{equation}

\subsubsection{Tiempo frente a comunicación}

\begin{figure}[H]
\centering
\fbox{
\parbox[c][6cm][c]{0.90\textwidth}{
\centering
\todoresult{Insertar gráfica de dispersión: \\
eje X = tiempo de autenticación; \\
eje Y = bytes transmitidos; \\
tamaño del punto = RAM pico.}
}}
\caption{Relación entre tiempo de autenticación, volumen de comunicación y memoria.}
\label{fig:time-bandwidth}
\end{figure}


% ============================================================
\subsubsection{Seguridad frente a rendimiento}

\begin{figure}[H]
\centering
\fbox{
\parbox[c][6cm][c]{0.90\textwidth}{
\centering
\todoresult{Insertar gráfica: \\
eje X = seguridad estimada; \\
eje Y = tiempo de autenticación.}
}}
\caption{Relación entre seguridad estimada y costo computacional.}
\label{fig:security-performance}
\end{figure}


\subsection{Discusión}
\label{sec:benchmark-discussion}

Los experimentos muestran que no existe necesariamente una alternativa que optimice
simultáneamente todas las dimensiones estudiadas.

Desde el punto de vista de rendimiento,
\todoresult{hallazgo}.

Respecto al consumo de memoria,
\todoresult{hallazgo}.

En términos de comunicación,
\todoresult{hallazgo}.

Respecto a la seguridad,
\todoresult{hallazgo}.

En términos de complejidad de implementación y despliegue,
\todoresult{hallazgo}.

Estos resultados evidencian un compromiso entre
\todoresult{trade-offs principales}.


% ============================================================
\subsubsection{Comparación con ECDSA}
\label{subsec:comparison-ecdsa}

ECDSA constituye el punto de referencia utilizado durante la evaluación.

Frente a esta alternativa, las construcciones basadas en retículos presentan
\todoresult{ventajas observadas}, mientras que presentan costos adicionales asociados
con \todoresult{desventajas observadas}.

En particular,
\todoresult{resultado cuantitativo importante}.


% ============================================================
\subsubsection{Comparación entre las variantes basadas en LWE}
\label{subsec:lwe-comparison}

Las variantes evaluadas muestran diferencias significativas en
\todoresult{métricas relevantes}.

LWE estándar presentó
\todoresult{resultado}.

Binary-LWE presentó
\todoresult{resultado}.

Ring-LWE presentó
\todoresult{resultado}.

LWR presentó
\todoresult{resultado}.

Estas diferencias permiten observar cómo modificaciones en la representación del
secreto, la estructura algebraica o el mecanismo utilizado para introducir
incertidumbre afectan los costos de una implementación real.


% ============================================================
\subsubsection{Evaluación del protocolo propuesto}
\label{subsec:proposed-evaluation}

El protocolo propuesto debe analizarse utilizando exactamente los mismos criterios que
las demás alternativas.

Los resultados muestran que
\todoresult{hallazgo principal}.

Su mayor ventaja observada corresponde a
\todoresult{ventaja}.

Su principal costo o limitación corresponde a
\todoresult{limitación}.

Frente a las demás variantes LWE,
\todoresult{comparación}.

Frente a ECDSA,
\todoresult{comparación}.

Por tanto, los resultados sugieren que el protocolo
\todoresult{interpretación prudente basada exclusivamente en los resultados obtenidos}.


% ============================================================
\subsection{Amenazas a la validez}
\label{sec:threats-validity}

Los resultados presentados deben interpretarse considerando las siguientes amenazas a
la validez.

\subsubsection{Validez interna}

Los tiempos de ejecución pueden verse afectados por procesos del sistema operativo,
planificación del procesador, cachés y otros factores externos. Para reducir este
efecto se utilizaron múltiples repeticiones, iteraciones de calentamiento y un entorno
experimental controlado.

\todoresult{Agregar otras amenazas internas identificadas.}


\subsubsection{Validez externa}

Los resultados corresponden al hardware, sistema operativo, versiones de software y
parámetros descritos en la Sección~\ref{sec:experimental-environment}. Por esta razón,
los valores absolutos pueden variar al utilizar hardware diferente.

Sin embargo, el diseño experimental busca permitir la reproducción del procedimiento
en otros entornos.

\todoresult{Agregar otras amenazas externas.}


\subsubsection{Validez de constructo}

Las métricas utilizadas representan diferentes dimensiones de eficiencia y
desplegabilidad, pero ninguna métrica individual captura completamente la calidad de
un protocolo criptográfico.

En particular, indicadores como líneas de código, cantidad de dependencias o número de
parámetros deben interpretarse únicamente como aproximaciones a la complejidad de
implementación y no como medidas absolutas de dificultad.


\subsubsection{Limitaciones de la estimación de seguridad}

Las estimaciones de seguridad dependen de los modelos de ataque, parámetros y
herramientas disponibles al momento de realizar el estudio.

\todoresult{Describir herramienta de estimación utilizada, versión, supuestos y
limitaciones.}


% ============================================================
\subsection{Reproducibilidad}
\label{sec:reproducibility}

Con el propósito de facilitar la reproducción de los experimentos, el repositorio del
proyecto conserva:

\begin{itemize}
    \item las implementaciones evaluadas;
    \item los parámetros criptográficos utilizados;
    \item los scripts de benchmarking;
    \item la configuración del entorno de ejecución;
    \item las versiones de las dependencias;
    \item las migraciones o estructura de almacenamiento;
    \item los resultados crudos de cada ejecución;
    \item los procedimientos utilizados para generar las métricas agregadas.
\end{itemize}

El experimento correspondiente a este documento fue ejecutado utilizando el commit:

\begin{center}
    \texttt{\todoresult{hash del commit}}
\end{center}

y la configuración experimental:

\begin{center}
    \texttt{\todoresult{identificador de configuración}}
\end{center}

